% Anhang: Referenz des hp-FEM-Plugins
% Quellen: python/plugins/hpfem/hpfem_plugin.py (Parameter, Aktionen), hpfem_job.py, hpfem_results.py,
% hpfem_map.py, hpfem_block.py, hpfem_cad.py, plugin.json; Blocktyp aus dem Paket rtt-fusion
% (rtt_fusion/graph/hpfem_grating.py)
\chapter{Referenz: hp-FEM-Plugin}\label{kap:plugin-hpfem}
\index{hp-FEM}

Dieser Anhang listet alle Aktionen, Parameter, Ergebnisse und Dateien des Plugins
\plugin{hp-FEM}. Die Bedienung erklärt \autoref{kap:hpfem}.

\section{Steckbrief}\label{sec:hp-ref-steckbrief}

\begin{referenz}{>{\raggedright\arraybackslash}p{3.4cm}>{\raggedright\arraybackslash}X}
\toprule
Eigenschaft & Wert \\
\midrule
\endhead
Name & \plugin{hp-FEM} \\
Kennung (\param{id}) & \code{hpfem} \\
Version & 0.4.0 \\
Plugin-API & 1 \\
Einstiegspunkt & \code{hpfem\_plugin:HpFemPlugin} \\
Python-Pakete (\param{requires}) & \code{numpy}, \code{hpfem>=0.5} \\
Ordner & \datei{python/plugins/hpfem} (wird mit der Anwendung installiert) \\
Blocktypen & \block{hpfem.grating} (mit \code{rtt-fusion} ab 0.5) \\
\bottomrule
\end{referenz}

\section{Python-Anforderungen}\label{sec:hp-ref-python}

\begin{referenz}{>{\raggedright\arraybackslash}p{4.2cm}>{\raggedright\arraybackslash}X}
\toprule
Paket & Wofür \\
\midrule
\endhead
\code{numpy} & Ergebnisse, Feldkarten, \gui{Show cell} \\
\code{hpfem} $\geq$ 0.5 & \gui{Check}, \gui{Solve}, \gui{Field map}, \gui{Library materials}; ohne das
  Paket meldet jede dieser Aktionen, dass es fehlt \\
\code{gmsh} & das Netz \code{gmsh} \\
\code{nurbforge} $\geq$ 0.7 & \gui{Shapes from CAD bodies} (ebener Schnitt geschlossener Körper) \\
\code{rtt\_fusion.\allowbreak graph.\allowbreak hpfem\_grating} (\code{rtt-fusion} $\geq$ 0.5) &
  Bereitstellen und Auswerten des Blocks \block{hpfem.grating} \\
\bottomrule
\end{referenz}

Ohne \code{hpfem} funktionieren: \gui{Load job}, \gui{Save job}, \gui{Show cell},
\gui{Shapes from CAD bodies} (der Name eines Körpers muss dann ein Material des Panels sein)
sowie die beiden Blockaktionen.

% ---------------------------------------------------------------------------------------------
\section{Aktionen}\label{sec:hp-ref-aktionen}

Die Aktionen stehen im Menü \menu{Plugins,hp-FEM}, als Knöpfe im Panel \gui{Plugin} und in der
Werkzeugspalte. Die beiden Blockaktionen stehen nur im Kontextmenü eines Blocks
\block{hpfem.grating} im Node-Editor.

\begin{referenz}{>{\raggedright\arraybackslash}p{3.6cm}>{\raggedright\arraybackslash}X>{\raggedright\arraybackslash}p{4.2cm}}
\toprule
Aktion & Wirkung & Liest / schreibt \\
\midrule
\endhead
\gui{Check} & baut das Netz und schätzt den Speicher einer Faktorisierung, ohne zu lösen &
  Protokoll (Netz, Schätzung, Warnung über 70\,\% RAM); Statusleiste \\
\gui{Field map} & rechnet einen Punkt des Sweeps und zeigt $\vec E$, $\vec H$ oder $\vec S$ auf
  einem Raster & Bild im Fenster \gui{Image}; \gui{Also write .npz} \\
\gui{Library materials} & listet die Materialien von hp-FEMs Bibliothek & Protokoll \\
\gui{Load job} & liest einen Job (.json) ins Panel & \gui{Job} \\
\gui{Save job} & schreibt den Job des Panels als .json & \gui{Save as}, sonst \gui{Job} \\
\gui{Shapes from CAD bodies} & schneidet die Körper in der Ebene CAD $y$ = \gui{Section at CAD y}
  und macht die Schnitte zu Polygonen mit Tags & \gui{Bodies}; Parameter \param{shapes},
  \param{materials}; Protokoll \\
\gui{Show cell} & zeichnet die Tags der Zelle auf dem Raster der Feldkarte, ohne zu rechnen & Bild
  \enquote{Tags}; Materialien im Protokoll \\
\gui{Solve} & rechnet den Job: Effizienzen je Ordnung, bei einem Sweep Kurven & Tabelle
  \enquote{Orders} und Balken bzw.\ Kurven und Tabelle \enquote{Points}; \datei{results.json} \\
\midrule
\gui{Store job in block} (Blockaktion) & der Job des Panels wird der des Blocks, ohne Einfall,
  Sweep und Karten & Blockparameter \param{job\_data}, \param{job}; Felder des Blockdialogs \\
\gui{Load block job into panel} (Blockaktion) & zeigt den Job des Blocks im Panel; Einfall und
  Sweep bleiben & Panel; \gui{Job} wird geleert \\
\bottomrule
\end{referenz}

% ---------------------------------------------------------------------------------------------
\section{Parameter}\label{sec:hp-ref-parameter}

Schlüssel (so steht er im Dokument und in \code{ctx.params}), Beschriftung, Typ, Standardwert und
Bedeutung, in der Reihenfolge des Panels. Längen stehen in der \gui{Unit} des Jobs
(\enquote{Einheit}). Die Typen sind die des FourierGrating-Anhangs
(\autoref{sec:fg-ref-parameter}).

\subsection{Ohne Gruppe: Formen}

\begin{referenz}{>{\raggedright\arraybackslash}p{3.1cm}>{\raggedright\arraybackslash}p{3.2cm}>{\raggedright\arraybackslash}X}
\toprule
Schlüssel / Beschriftung & Typ / Standard & Bedeutung \\
\midrule
\endhead
\param{shapes}\newline\gui{Shapes (the last on top)} & Liste; Standard: ein Rechteck, Tag 3,
  $x = -100$, $y = 0$, 200 × 148 & Formen der Zelle; die letzte liegt oben \\
\midrule
\multicolumn{3}{@{}l}{\textit{Felder eines Eintrags}} \\
\gui{Type} & Auswahl: \code{rectangle}, \code{trapezoid}, \code{ellipse}, \code{polygon},
  \code{json} & Form; welche Felder gelten, steht in \autoref{sec:hp-zelle} \\
\gui{Tag} & Ganzzahl, 3 (0 … 10000) & Tag der Form, ihr Material in \gui{Materials by tag} \\
\gui{x}, \gui{y} & Zahl, 0 & Ecke (Rechteck), Mitte/Basis (Trapez), Mitte (Ellipse) \\
\gui{width}, \gui{height} & Zahl, 0 & Breite, Höhe (Rechteck; Höhe auch Trapez) \\
\gui{bottom}, \gui{top} & Zahl, 0 & untere und obere Breite (Trapez) \\
\gui{rx}, \gui{ry} & Zahl, 0 & Halbachsen (Ellipse) \\
\gui{Points (polygon)} & Tabelle \gui{x}, \gui{y} & Ecken gegen den Uhrzeigersinn \\
\gui{Shape as JSON} & JSON & jede andere Form des Job-Formats (Typ \code{json}) \\
\bottomrule
\end{referenz}

\subsection{Gruppen File und Structure}

\begin{referenz}{>{\raggedright\arraybackslash}p{3.1cm}>{\raggedright\arraybackslash}p{3.2cm}>{\raggedright\arraybackslash}X}
\toprule
Schlüssel / Beschriftung & Typ / Standard & Bedeutung \\
\midrule
\endhead
\param{job\_file}\newline\gui{Job} & Datei & Job für \gui{Load job}; eine Netzdatei gilt relativ
  zu ihm \\
\param{save\_as}\newline\gui{Save as} & Datei (speichern) & Ziel von \gui{Save job} \\
\param{results\_file}\newline\gui{Results in folder of} & Datei (speichern) & in diesen Ordner
  schreibt \gui{Solve} \datei{results.json}; leer: ein temporärer Ordner \\
\param{name}\newline\gui{Name} & Text, \code{bar on SiO2} & Name des Jobs (Titel der Ergebnisse) \\
\param{unit}\newline\gui{Unit} & Auswahl \code{nm}, \code{µm}, \code{mm}, \code{m}; \code{nm} &
  Einheit aller Längen und der Wellenlänge \\
\param{period}\newline\gui{Period} & Zahl $\geq 0$, 400 & Periode in $x$ \\
\param{y\_bottom}\newline\gui{Cell bottom y} & Zahl, $-800$ & Unterkante der Zelle (PML innen) \\
\param{y\_top}\newline\gui{Cell top y} & Zahl, 950 & Oberkante der Zelle \\
\param{background\_tag}\newline\gui{Background tag} & Ganzzahl, 1 & Tag, wo keine Schicht und keine
  Form liegt \\
\param{slabs}\newline\gui{Slabs (full period)} & Tabelle \gui{tag}, \gui{y0}, \gui{y1}; Tag 2 von
  $-800$ bis 0 & Schichten über die ganze Periode \\
\bottomrule
\end{referenz}

\subsection{Gruppen Materials und Stack}

\begin{referenz}{>{\raggedright\arraybackslash}p{3.1cm}>{\raggedright\arraybackslash}p{3.2cm}>{\raggedright\arraybackslash}X}
\toprule
Schlüssel / Beschriftung & Typ / Standard & Bedeutung \\
\midrule
\endhead
\param{materials}\newline\gui{Materials by tag} & Tabelle \gui{tag}, \gui{material}; 1 \code{air},
  2 \code{SiO2}, 3 \code{eps=2.25+0i} & Material je Tag (Schreibweisen:
  \autoref{sec:hp-materialien}) \\
\param{incidence\_medium}\newline\gui{Incidence medium (above)} & Text, \code{air} & Medium über
  dem Stapel (verlustfrei) \\
\param{substrate}\newline\gui{Substrate (below)} & Text, \code{SiO2} & Medium unter dem Stapel \\
\param{stack\_layers}\newline\gui{Background layers (top down)} & Tabelle \gui{material},
  \gui{thickness}; leer & Schichten des Hintergrundstapels von oben nach unten \\
\param{stack\_top}\newline\gui{Top interface y} & Zahl, 0 & Lage der obersten Grenzfläche \\
\bottomrule
\end{referenz}

\subsection{Gruppe Incidence}

\begin{referenz}{>{\raggedright\arraybackslash}p{3.1cm}>{\raggedright\arraybackslash}p{3.2cm}>{\raggedright\arraybackslash}X}
\toprule
Schlüssel / Beschriftung & Typ / Standard & Bedeutung \\
\midrule
\endhead
\param{polarisation}\newline\gui{Polarisation} & Auswahl \code{s}, \code{p}; \code{p} & \code{s}:
  $E_z$ bei $\varphi = 0$ (TE), \code{p}: $H_z$ (TM) \\
\param{wavelength}\newline\gui{Wavelength} & Text, \code{400} & Zahl, Liste oder
  \code{Start:Stopp:Anzahl} (Einheit) \\
\param{theta\_deg}\newline\gui{$\theta$ / deg} & Text, \code{50} & Polarwinkel; Zahl, Liste oder
  Bereich (nur einer von Wellenlänge und $\theta$ als Sweep) \\
\param{phi\_deg}\newline\gui{$\varphi$ / deg} & Zahl, 30 & Azimut; $\neq 0$: konischer Einfall \\
\bottomrule
\end{referenz}

\subsection{Gruppen Mesh und Solver}

\begin{referenz}{>{\raggedright\arraybackslash}p{3.1cm}>{\raggedright\arraybackslash}p{3.2cm}>{\raggedright\arraybackslash}X}
\toprule
Schlüssel / Beschriftung & Typ / Standard & Bedeutung \\
\midrule
\endhead
\param{mesh\_kind}\newline\gui{Mesh} & Auswahl \code{structured}, \code{gmsh}, \code{file};
  \code{structured} & Art des Netzes (\autoref{sec:hp-netz}) \\
\param{mesh\_nx}\newline\gui{Columns (structured)} & Ganzzahl, 16 & Spalten über die Periode \\
\param{mesh\_rows}\newline\gui{Row bands (structured, bottom up)} & Tabelle \gui{y0}, \gui{y1},
  \gui{rows}; $-800$…0: 32, 0…148: 6, 148…950: 32 & Zeilenbänder \\
\param{gmsh\_interface\_factor}\newline\gui{Interface factor (gmsh)} & Zahl 0,01 … 10; 0,5 &
  Verfeinerung an Grenzflächen \\
\param{mesh\_file}\newline\gui{Mesh file (file)} & Datei & periodisches Gmsh-Netz (.msh) \\
\param{mesh\_scale}\newline\gui{Scale to metres (file)} & Zahl, 1 & Faktor der Netzkoordinaten auf
  Meter \\
\param{order}\newline\gui{Polynomial order p} & Ganzzahl 1 … 12; 3 & Ordnung der Elemente \\
\param{pml\_top}, \param{pml\_bottom}\newline\gui{PML top}, \gui{PML bottom} & Zahl $\geq 0$; 419
  & Dicke der PML innerhalb der Zelle; 0: hp-FEMs Empfehlung \\
\param{bottom}\newline\gui{Bottom} & Auswahl \code{pml}, \code{pec}, \code{modal}; \code{pml} &
  Abschluss unten \\
\param{orders\_max}\newline\gui{Orders up to |m|} & Ganzzahl 0 … 200; 2 & ausgewertete Ordnungen \\
\param{check}\newline\gui{Diagnostics before solving} & Schalter, ein & hp-FEMs Prüfungen \\
\bottomrule
\end{referenz}

\subsection{Gruppen Results und Field map}

\begin{referenz}{>{\raggedright\arraybackslash}p{3.1cm}>{\raggedright\arraybackslash}p{3.2cm}>{\raggedright\arraybackslash}X}
\toprule
Schlüssel / Beschriftung & Typ / Standard & Bedeutung \\
\midrule
\endhead
\param{min\_efficiency}\newline\gui{Orders from efficiency} & Zahl 0 … 1; 0,01 & ein Sweep zeigt
  eine Ordnung, die das irgendwo erreicht \\
\param{evanescent}\newline\gui{List evanescent orders} & Schalter, aus & abklingende Ordnungen in
  der Tabelle \\
\param{map\_quantity}\newline\gui{Field} & Auswahl \code{E}, \code{H}, \code{S}; \code{E} & Feld
  der Karte (SI) \\
\param{map\_show}\newline\gui{Show} & Auswahl $|F|^2$, $|F_x|$, $|F_y|$, $|F_z|$, Re/Im $F_x$,
  $F_y$, $F_z$; $|F|^2$ & gezeigte Größe \\
\param{map\_point}\newline\gui{Point of the sweep} & Ganzzahl $\geq 1$; 1 & Punkt der Karte \\
\param{map\_scattered}\newline\gui{Scattered field only} & Schalter, aus & nur das Streufeld \\
\param{map\_periods}\newline\gui{Periods} & Ganzzahl 1 … 50; 2 & Perioden in $x$ (wenn \gui{x
  from} = \gui{x to}) \\
\param{map\_x\_from}, \param{map\_x\_to} & Zahl, 0 & Bereich in $x$ \\
\param{map\_nx}\newline\gui{Points x} & Ganzzahl 2 … 4000; 201 & Punkte in $x$ \\
\param{map\_y\_from}, \param{map\_y\_to} & Zahl, 0 & Bereich in $y$; gleich: Zelle ohne PML \\
\param{map\_ny}\newline\gui{Points y} & Ganzzahl 2 … 4000; 301 & Punkte in $y$ \\
\param{map\_npz}\newline\gui{Also write .npz} & Datei (speichern) & Karte als NumPy-Archiv \\
\bottomrule
\end{referenz}

\subsection{Gruppen CAD cell und Advanced}

\begin{referenz}{>{\raggedright\arraybackslash}p{3.1cm}>{\raggedright\arraybackslash}p{3.2cm}>{\raggedright\arraybackslash}X}
\toprule
Schlüssel / Beschriftung & Typ / Standard & Bedeutung \\
\midrule
\endhead
\param{cad\_bodies}\newline\gui{Bodies} & Körper (Schalen, mehrere) & geschlossene Körper einer
  Elementarzelle \\
\param{cad\_unit}\newline\gui{Bodies drawn in} & Auswahl \code{nm}, \code{µm}, \code{mm};
  \code{µm} & was eine CAD-Längeneinheit ist \\
\param{cad\_section\_y}\newline\gui{Section at CAD y} & Zahl, 0 & Lage der Schnittebene (CAD) \\
\param{cad\_points}\newline\gui{Points per curved edge} & Ganzzahl 4 … 2000; 48 & Punkte je
  gekrümmter Kante \\
\param{cad\_mode}\newline\gui{Shapes} & Auswahl \code{replace shapes}, \code{add shapes} & Formen
  ersetzen oder ergänzen \\
\param{advanced}\newline\gui{Further keys (JSON)} & JSON, leer & alle übrigen Schlüssel des Jobs,
  beim Speichern eingefügt \\
\bottomrule
\end{referenz}

% ---------------------------------------------------------------------------------------------
\section{Panel und Job-Datei}\label{sec:hp-ref-zuordnung}

\begin{referenz}{>{\raggedright\arraybackslash}p{5.0cm}>{\raggedright\arraybackslash}X}
\toprule
Panel & Job (Schema-Version 1) \\
\midrule
\endhead
\param{name}, \param{unit} & \code{name}, \code{unit} (Faktor auf Meter) \\
Gruppe \gui{Structure}, \param{shapes} & \code{model.period}, \code{y\_bottom}, \code{y\_top},
  \code{background\_tag}, \code{slabs}, \code{shapes} \\
\param{materials} & \code{materials.<Tag>}: Name, \code{\{\textquotedbl{}n\textquotedbl{}: …\}} oder \code{\{\textquotedbl{}eps\textquotedbl{}: …\}} \\
Gruppe \gui{Stack} & \code{stack.incidence}, \code{substrate}, \code{layers}, \code{top} \\
Gruppe \gui{Incidence} & \code{incidence}; ein Sweep als \code{sweep.wavelength} bzw.\
  \code{sweep.theta\_deg} (Liste oder \code{\{\textquotedbl{}start\textquotedbl{}, \textquotedbl{}stop\textquotedbl{}, \textquotedbl{}count\textquotedbl{}\}}) \\
Gruppe \gui{Mesh} & \code{mesh.structured} (\code{nx}, \code{rows}), \code{mesh.gmsh}
  (\code{interface\_factor}) oder \code{mesh.file} (\code{file}, \code{scale}) \\
Gruppe \gui{Solver} & \code{solver.order}, \code{pml} (nur Dicken $> 0$), \code{bottom},
  \code{orders\_max}, \code{check} \\
\param{advanced} & alle übrigen Schlüssel (z.\,B.\ \code{maps}, \code{model.x0}), unverändert \\
\bottomrule
\end{referenz}

% ---------------------------------------------------------------------------------------------
\section{Ergebnisse}\label{sec:hp-ref-ergebnisse}

\begin{referenz}{>{\raggedright\arraybackslash}p{3.2cm}>{\raggedright\arraybackslash}p{3.6cm}>{\raggedright\arraybackslash}X}
\toprule
Ergebnis & Art & Inhalt \\
\midrule
\endhead
\enquote{Orders} & Tabelle (ein Punkt) & \code{side} (R, T), \code{m}, \code{efficiency},
  \code{propagating} \\
\enquote{Efficiency per order} & Balken (ein Punkt) & je ausbreitungsfähiger Ordnung
  \enquote{R~$m$} bzw.\ \enquote{T~$m$} \\
$\Sigma R$, $\Sigma T$, $A$, \enquote{R~$m$}, \enquote{T~$m$} & Kurven (Sweep) & über $\lambda$ in
  der Einheit des Jobs oder über $\theta$ in Grad \\
\enquote{Points} & Tabelle (Sweep) & \code{point}, $\lambda$, \code{$\theta$ / deg}, \code{R}, \code{T},
  \code{A}, \code{balance}, \code{DoFs}, \code{seconds} \\
$|E|^2$, \enquote{Re Sy}, … & Bild (\gui{Field map}) & die gezeigte Größe; Achsen in der Einheit
  des Jobs; leere Punkte außerhalb der Zelle \\
\enquote{Tags} & Bild (\gui{Show cell}) & Tag je Rasterpunkt \\
\bottomrule
\end{referenz}

Dateien: \datei{results.json} des Runners (alle Punkte mit Ordnungen, Amplituden, Diagnosen,
Zeiten) im Ordner von \gui{Results in folder of}; Feldkarten als \datei{maps\_<Punkt>\_<Karte>.npz}
(\code{x}, \code{y} in m, \code{values} komplex $(n_y, n_x, 3)$, \code{quantity},
\code{wavelength}).

% ---------------------------------------------------------------------------------------------
\section{Workflow-Block}\label{sec:hp-ref-block}

Der Block \block{hpfem.grating} mit allen Parametern steht in \autoref{sec:block-hpfem-grating}.
Der Dialog eines Blocks mit Job zeigt dazu Felder des Jobs (JSON-Zeiger in \param{job\_data}):

\begin{referenz}{>{\raggedright\arraybackslash}p{4.4cm}>{\raggedright\arraybackslash}p{2.4cm}>{\raggedright\arraybackslash}X}
\toprule
Feld & Gruppe & Zeiger \\
\midrule
\endhead
\gui{Period} & \gui{Grating} & \code{/job\_data/model/period} \\
\gui{Polynomial order} & \gui{Solver} & \code{/job\_data/solver/order} \\
\gui{Tag $t$: n} bzw.\ \gui{Tag $t$: n, k} & \gui{Materials} & \code{/job\_data/materials/<t>/n} (nur
  Materialien, die als $n$ gegeben sind) \\
\bottomrule
\end{referenz}
